\documentclass[UTF8,11pt]{ctexart}

\usepackage[a4paper,top=2.25cm,bottom=2.2cm,left=2.4cm,right=2.4cm,headheight=15pt]{geometry}
\usepackage{fontspec}
\setmainfont{TeX Gyre Pagella}
\setsansfont{TeX Gyre Heros}
\setmonofont{DejaVu Sans Mono}[Scale=0.84]
\setCJKmainfont{Noto Serif CJK SC}
\setCJKsansfont{Noto Sans CJK SC}
\setCJKmonofont{Noto Sans Mono CJK SC}

\usepackage{amsmath,amssymb,bm}
\usepackage{graphicx}
\usepackage{booktabs,tabularx,array,longtable}
\usepackage{siunitx}
\usepackage[dvipsnames]{xcolor}
\usepackage{tikz}
\usepackage[most]{tcolorbox}
\usepackage{enumitem}
\usepackage{caption}
\usepackage{float}
\usepackage{fancyhdr}
\usepackage{titlesec}
\usepackage{listings}
\usepackage{xurl}
\usepackage{hyperref}
\usepackage{lastpage}

\definecolor{Navy}{HTML}{153F66}
\definecolor{Blue}{HTML}{2F6F9F}
\definecolor{Teal}{HTML}{2A8C82}
\definecolor{Green}{HTML}{4D8F62}
\definecolor{Orange}{HTML}{D9822B}
\definecolor{LightBlue}{HTML}{EEF5FA}
\definecolor{LightGray}{HTML}{F4F6F8}
\definecolor{Dark}{HTML}{25313B}

\hypersetup{
  colorlinks=true,
  linkcolor=Navy,
  urlcolor=Blue,
  citecolor=Teal,
  pdftitle={多月份真实环境锚定的深海声传播损失快速计算技术报告},
  pdfauthor={Ocean Acoustic Surrogate Project},
  pdfsubject={GEBCO、WOA23、Bellhop 与改进型 Fourier Neural Operator}
}

\graphicspath{{assets/}{../assets/}}
\captionsetup{font=small,labelfont={bf,color=Navy}}
\setlist{leftmargin=2em,itemsep=0.2em,topsep=0.3em}
\renewcommand{\arraystretch}{1.22}
\setlength{\tabcolsep}{5pt}
\setcounter{tocdepth}{2}
\setcounter{secnumdepth}{3}
\titleformat{\section}{\LARGE\bfseries\sffamily\color{Navy}}{\thesection}{0.7em}{}
  [\vspace{0.2em}\color{Blue}\titlerule]
\titleformat{\subsection}{\Large\bfseries\sffamily\color{Navy}}{\thesubsection}{0.7em}{}
\titleformat{\subsubsection}{\large\bfseries\sffamily\color{Blue}}{\thesubsubsection}{0.6em}{}

\pagestyle{fancy}
\fancyhf{}
\fancyhead[L]{\small\sffamily\color{Navy}深海声传播损失快速计算}
\fancyhead[R]{\small\sffamily\color{Navy}项目技术报告 V1.5}
\fancyfoot[L]{\small\color{gray}Ocean Acoustic Surrogate Project}
\fancyfoot[R]{\small\color{gray}\thepage\,/\,\pageref{LastPage}}
\renewcommand{\headrulewidth}{0.45pt}

\newcolumntype{Y}{>{\raggedright\arraybackslash}X}
\newcolumntype{L}[1]{>{\raggedright\arraybackslash}p{#1}}
\newtcolorbox{summarybox}{enhanced,breakable,colback=LightBlue,colframe=Blue,
  boxrule=0.7pt,arc=1mm,left=4mm,right=4mm,top=3mm,bottom=3mm}
\lstdefinestyle{repro}{basicstyle=\ttfamily\footnotesize,backgroundcolor=\color{LightGray},
  frame=single,rulecolor=\color{gray!35},breaklines=true,columns=fullflexible,
  keepspaces=true,showstringspaces=false,xleftmargin=0.8em,framexleftmargin=0.6em}
\newcommand{\figfull}[3]{\begin{figure}[H]\centering
  \includegraphics[width=#1\textwidth]{#2}\caption{#3}\end{figure}}
\newcommand{\TL}{\mathrm{TL}}
\newcommand{\RMSE}{\mathrm{RMSE}}

% Frozen-result macros; replaced from sealed artifacts before release.
\newcommand{\DatasetHash}{\texttt{9c90e051851e2408\ldots 0fda0a61}}
\newcommand{\BestModel}{SeaBAR-FNO}
\newcommand{\BestRMSE}{0.702}
\newcommand{\BestMAE}{0.295}
\newcommand{\BestPninety}{1.031}
\newcommand{\BestWorst}{1.320}
\newcommand{\GlobalBaseline}{4.954}
\newcommand{\RelativeReduction}{85.83}
\newcommand{\GpuPninetyfive}{4.70}
\newcommand{\CpuPninetyfive}{100.58}
\newcommand{\BellhopMedian}{8.10}

\begin{document}

\begin{titlepage}
\begin{tikzpicture}[remember picture,overlay]
  \fill[Navy] (current page.north west) rectangle ([yshift=-1.4cm]current page.north east);
  \fill[Blue] ([yshift=-1.4cm]current page.north west) rectangle
    ([yshift=-1.52cm]current page.north east);
\end{tikzpicture}
\vspace*{2.0cm}
{\sffamily\color{Blue}\Large 项目技术报告}\par
\vspace{0.75cm}
{\sffamily\bfseries\color{Navy}\fontsize{26}{34}\selectfont
多月份真实环境数据锚定的\par
深海声传播损失快速计算}\par
\vspace{0.65cm}
{\sffamily\color{Dark}\Large
GEBCO 地形、WOA23 月度声速剖面与 SeaBAR-FNO}\par

\vspace{1.45cm}
\begin{summarybox}
\centering
\begin{tabular}{L{3.4cm}L{8.0cm}}
\textbf{核心任务} & 从声速剖面与距离相关地形直接预测二维传播损失场 \\
\textbf{固定参数} & 1 kHz、50 km、50 m 声源深度、96$\times$256 输出网格 \\
\textbf{环境数据} & GEBCO 2026 巴士海峡候选区；WOA23 3/6/12 月温盐经 TEOS-10 转声速 \\
\textbf{参考标签} & Bellhop 非相干传播损失，25,600 条射线/场 \\
\textbf{主要结果} & 测试 RMSE \BestRMSE{} dB；均值场 \GlobalBaseline{} dB；
  A100 P95 \GpuPninetyfive{} ms \\
\textbf{报告版本} & V1.5，2026-08-28 \\
\end{tabular}
\end{summarybox}
\vfill
{\sffamily\color{Navy}\Large Ocean Acoustic Surrogate Project}\par
\vspace{0.3cm}
{\color{Dark}面向固定典型深海任务域的高质量仿真与快速代理计算研究}\par
\vspace{1.0cm}
\end{titlepage}

\pagenumbering{roman}
\section*{摘\quad 要}
\addcontentsline{toc}{section}{摘要}

本项目面向频率 1 kHz、最大传播距离 50 km、声源深度 50 m 的典型深海声传播损失快速
计算任务，研究从声速剖面与距离相关海底地形到二维非相干传播损失（Transmission Loss,
TL）场的代理映射。项目的核心验收要求为：相对于 Bellhop 数值参考，独立测试集平均误差
不超过 2 dB；单场模型计算时间不超过 100 ms。

为兼顾真实环境依据与首阶段可控性，数据集采用“\textbf{少量代表地形、多月份真实 SSP}”
的受控复杂度设计。海底形态取自 GEBCO 2026 巴士海峡候选区子集：从候选中心沿 24 个
方位提取 50 km 剖面，经 9 km 平滑、深水包络归一化和低维重采样后，预先选取四条声场
差异较大的代表剖面。声速基准来自同一区域 WOA23 0.25$^\circ$ 的 12 个月温盐气候态，
通过 TEOS-10 转换为声速后，以受约束三中心分析选出 3、6、12 月代表季节变化，并叠加
幅度受限的四参数平滑扰动。正式数据集包含 384 个环境样本，按地形和月份联合分层划分为
288/48/48 个训练、验证和独立测试样本。

参考标签由 Bellhop 非相干场生成。覆盖三个月份与四类地形的 8 场新审计表明，12,800 到
25,600 射线的平均差为 0.207 dB、最差样本为 0.302 dB，因此正式数据固定使用 25,600
条射线。方法上提出 Seasonal and Bathymetry-Aware Residual Fourier Neural Operator
（SeaBAR-FNO）：通过双环境上下文编码、各向异性谱传播核、全局--局部残差融合、非周期
边界保护及训练域残差解码，学习 SSP 与地形共同决定的传播结构。

密封测试结果显示，最终模型 \BestModel{} 的 RMSE 为 \BestRMSE{} dB、MAE 为
\BestMAE{} dB、P90 单样本 RMSE 为 \BestPninety{} dB，均满足 2 dB 指标。相比训练集全局
均值场 \GlobalBaseline{} dB，RMSE 相对下降 \RelativeReduction{}\%。A100 GPU 上 batch=1
完整推理 P95 为 \GpuPninetyfive{} ms，满足 100 ms 指标。三个月份分组 RMSE 均低于
0.75 dB，四类地形分组均低于 0.87 dB。本报告中的“参考场”均指公开
环境数据驱动的 Bellhop 高质量仿真结果，不等同于海试真值。

\vspace{0.5cm}
\noindent\textbf{关键词：}水声传播；传播损失；Bellhop；GEBCO；WOA23；
Fourier Neural Operator；代理模型

\clearpage
\tableofcontents
\clearpage
\pagenumbering{arabic}

\section{任务理解与技术目标}

\subsection{核心任务}

给定一维声速剖面 $c(z)$、距离相关海底深度 $b(r)$ 以及固定的频率、声源和海底介质参数，
项目需要构造代理算子
\begin{equation}
  \mathcal{G}_\theta:\big(c(z),b(r)\big)\longmapsto \TL(z,r),
\end{equation}
一次性输出完整的深度--距离二维 TL 场。与逐个接收点回归不同，神经算子直接学习函数到
函数的映射，并在统一网格上同时恢复传播扩散、声道折射、海底反射和会聚区结构。

\figfull{0.98}{fig01_project_pipeline.pdf}{项目技术路线：公开环境数据经可追溯处理形成受控任务域，
Bellhop 提供高质量参考，改进型 FNO 完成毫秒级预测。}

\subsection{固定场景与硬指标}

\begin{table}[H]
\centering
\caption{核心场景、数据网格与验收指标}
\label{tab:contract}
\begin{tabularx}{\textwidth}{L{3.5cm}L{4.1cm}Y}
\toprule
项目 & 冻结设置 & 说明 \\
\midrule
声源频率 & 1000 Hz & 单频 \\
声源深度 & 50 m & 每个样本固定 \\
最大传播距离 & 50 km & 0.1--50 km，共 256 点 \\
接收深度 & 10--1990 m & 共 96 点 \\
海底介质 & 流体沙质半空间 & 1700 m/s，2000 kg/m$^3$，0.8 dB/$\lambda$ \\
地形水深包络 & 2000--4800 m & 由 GEBCO 形态受控归一化 \\
参考场 & Bellhop 非相干 TL & 25,600 rays/field \\
精度门槛 & RMSE、MAE 均不超过 2 dB & 样本级密封测试 \\
时延门槛 & P95 不超过 100 ms & batch=1，含输入拷贝和反归一化 \\
任务难度门槛 & 全局均值场 RMSE 大于 4 dB & 防止简单数据造成无意义达标 \\
\bottomrule
\end{tabularx}
\end{table}

本阶段不追求跨海区、跨频率或任意底质的通用模型，而是在满足固定参数的前提下，先验证
真实数据锚定环境中的准确率和时延闭环。测试样本与训练样本共享四类地形定义，但 SSP
参数独立抽样；因此结果衡量的是受控任务域内的插值能力，而非未见地形类别的外推能力。

\section{环境数据与参考数据集}

\subsection{数据来源与可追溯性}

\begin{table}[H]
\centering
\caption{环境数据来源及其在本项目中的用途}
\begin{tabularx}{\textwidth}{L{3.1cm}L{4.2cm}Y}
\toprule
数据 & 冻结来源 & 处理与用途 \\
\midrule
海底地形 & GEBCO 2026 巴士海峡候选框，中心 20.5$^\circ$N、121.5$^\circ$E &
沿方位提取 50 km 地形剖面，构造 Bellhop 距离相关边界和模型地形通道 \\
温度、盐度 & WOA23 decadal 0.25$^\circ$ 月度气候态，3、6、12 月 &
取候选中心附近网格；1500 m 以下使用年度气候态补齐 \\
声速 & WOA 温盐经 TEOS-10/GSW 转换 &
形成 0--2000 m 三类代表性月度深海 SSP 基准 \\
参考声场 & Bellhop &
在冻结频率、源深、底质和网格上生成非相干 TL \\
\bottomrule
\end{tabularx}
\end{table}

GEBCO 区域原始文件 SHA-256 为
\path{1e4a2ed7c6742499d378f5f90c99b6cb3ce066d6ab034593473afbddb24ad2eb}。
WOA23 处理文件包含 12 个月温盐和转换声速，其 SHA-256 为
\path{38e548acc493dba35580735248581b4dced886b5c18aee2c5c5437d92d56db16}。
对 12 个月基准剖面逐点标准化后，以剖面间平方距离进行穷举三中心分析：不受约束最优组为
3/7/11 月；为保留前期六月数据锚点，约束候选包含 6 月后得到 3/6/12 月，其覆盖目标函数
仅比无约束最优高 3.1\%。该选择用三个真实月度型态表示季节差异，同时避免一次性引入全部
12 个月和跨海区变化。

\subsection{GEBCO 地形剖面的受控构造}

地形准备不是随机生成四条曲线，而是执行下列固定流程：
\begin{enumerate}
  \item 在候选中心沿 $0^\circ,15^\circ,\ldots,345^\circ$ 共 24 个方位，以 1 km 间隔
  提取 50 km GEBCO 水深序列；
  \item 使用 9 km 移动平均抑制单格噪声和过窄地形起伏；
  \item 保留每条剖面的相对形态，将水深仿射映射到 2000--4800 m 的任务深水包络；
  \item 以约 5 km 间隔降维，同时显式保留局部最浅点和最深点；
  \item 在名义 SSP 和 3200 射线预筛条件下，计算候选四剖面组合的均值场 RMSE，
  在正式样本生成前选定 $30^\circ,90^\circ,210^\circ,225^\circ$ 四个方位。
\end{enumerate}

预筛只使用一个名义 SSP 和低成本标签，不使用正式训练、验证或测试样本。24 方位中选定
四剖面的预筛均值场 RMSE 为 5.631 dB；其目的在于构造既小、又具有明确区分度的受控
任务域，而不是扩大数据覆盖范围。

\figfull{0.99}{fig02_real_environment.pdf}{巴士海峡候选区位置、四个 50 km 建模方向、WOA23
3/6/12 月 SSP 族与经筛选的低维 GEBCO 地形剖面。橙色星号为剖面起点。}

\subsection{WOA23 月度声速剖面与受控扰动}

WOA23 温盐转换使用 TEOS-10：由深度和纬度得到压力 $p$，将实用盐度 $S_P$ 和原位温度
$t$ 转换为绝对盐度 $S_A$、保守温度 $\Theta$，再计算
\begin{equation}
  c(z)=c_{\mathrm{TEOS10}}\big(S_A(z),\Theta(z),p(z)\big).
\end{equation}
三个月份基准剖面均在约 1000--1300 m 形成声速低值区并在深层缓慢回升，但上层结构具有
可辨季节差异：6 月表层约 1542.68 m/s，3 月和 12 月表层约 1535.7--1536.0 m/s。

为构造每个月份内部的平滑变化，每个样本采用四个可解释参数：全局偏移
$\delta_c\in[-0.55,0.55]$ m/s、温跃层幅度 $a_t\in[-0.8,0.8]$ m/s、声道轴平移
$\delta_z\in[-35,35]$ m 和深层梯度 $a_d\in[-0.55,0.55]$ m/s。对月份
$m_i\in\{3,6,12\}$，连续 SSP 写为
\begin{equation}
 c_i(z)=c_{0,m_i}(\operatorname{clip}(z-\delta_{z,i}))+\delta_{c,i}
 +a_{t,i}\exp\!\left[-\frac{(z-250)^2}{2(220)^2}\right]
 +a_{d,i}\operatorname{clip}\!\left(\frac{z-1000}{1000},0,1\right).
\end{equation}
四维参数采用固定种子的分块 Latin hypercube 采样。一个 96 场设计块完整覆盖
$4$ 类地形$\times3$ 个月份，每个组合包含 8 场；正式设计由四个块组成，每个环境组合
共 32 场。所得 SSP 总体范围为 1482.52--1543.34 m/s，样本族逐点跨度中位数为
2.23 m/s、最大值为 12.53 m/s。该参数化只产生平滑剖面，不制造网格尺度随机噪声。

\subsection{Bellhop 高质量标签}

每个环境样本生成一个 96$\times$256 的非相干 TL 场。海底采用声速 1700 m/s、密度
2000 kg/m$^3$、衰减 0.8 dB/$\lambda$ 的流体半空间。正式计算前，从三个月份和四类地形
中分层抽取 8 个独立环境，在 12,800 和 25,600 条射线下重新运行，并以场差评估数值稳定性。

\begin{table}[H]
\centering
\caption{射线数逐级收敛审计}
\begin{tabular}{lrrr}
\toprule
比较 & 平均 RMSE/dB & 最差 RMSE/dB & 平均 P95 绝对差/dB \\
\midrule
12,800 $\rightarrow$ 25,600 & 0.207 & 0.302 & 0.158 \\
\bottomrule
\end{tabular}
\end{table}

\figfull{0.94}{fig03_label_convergence.pdf}{射线数增加时的相邻级别 TL 差与单场计算成本。
正式标签固定使用 25,600 条射线。}

\subsection{样本规模、分层划分与文件契约}

正式数据集包含 384 个样本，即 4 类地形$\times$3 个月份$\times$32 个独立 SSP 扰动。
每个地形--月份组合固定划分为 24 个训练、4 个验证和 4 个测试样本，汇总为 288/48/48。
划分单位是完整环境样本，绝不把同一二维场的网格点随机分散到多个 split。测试标签在模型
结构和超参数确定后才用于一次性评分。扩容过程逐元素核验并复用同一季节设计的前 96 个
高射线标签，仅为新增设计点补算 288 个 Bellhop 场。

每个样本保存 TL、有效掩膜、SSP、地形、坐标、参数和 Bellhop case 元数据。打包文件
\path{dataset.npz} 的 SHA-256 为 \DatasetHash；\path{manifest.json} 同时记录配置哈希、
项目 Git SHA、声学求解器 Git SHA、Python 版本、生成耗时、覆盖率和失败清单。大型数据
位于外部数据盘，不写入 Git；Git 中只保留配置、代码、哈希、轻量指标和报告。

\section{SeaBAR-FNO：季节与地形感知的残差 Fourier 神经算子}

本项目将最终方法命名为 \textbf{SeaBAR-FNO}（Seasonal and Bathymetry-Aware Residual
Fourier Neural Operator）。名称对应任务中的两类可变环境量和一个核心学习策略：Seasonal
表示模型接收月度 SSP；Bathymetry-Aware 表示距离相关海底边界显式进入网络；Residual
表示模型预测训练均值场之上的变化，而不是从零拟合全部 TL。该名称是对实际计算图的概括，
并不额外引入无法复核的物理先验。

\subsection{符号、索引与离散维度}

为避免把“连续函数”“二维数组”和“神经网络张量”混为一谈，首先固定全文符号。样本下标
$i=1,\ldots,N$；SSP 原始深度点下标 $j=1,\ldots,N_s$；接收深度下标
$p=1,\ldots,N_z$；距离下标 $q=1,\ldots,N_r$；隐藏通道下标
$h=1,\ldots,H$。本轮最终设置为
\begin{equation}
  N=384,\qquad N_s=41,\qquad N_z=96,\qquad N_r=256,\qquad H=32.
\end{equation}

\begin{table}[H]
\centering
\caption{主要符号、物理意义和程序维度}
\begin{tabularx}{\textwidth}{L{2.2cm}L{3.2cm}L{3.0cm}Y}
\toprule
符号 & 物理意义 & 单样本维度 & 取值或单位 \\
\midrule
$c_i(z_j)$ & 第 $i$ 个 SSP & $[N_s]=[41]$ & m/s，深度 0--2000 m \\
$b_i(r_q)$ & 距离相关海底深度 & $[N_r]=[256]$ & m，正值表示向下 \\
$Y_i(p,q)$ & Bellhop 非相干 TL 标签 & $[N_z,N_r]=[96,256]$ & dB \\
$M_i(p,q)$ & 标签有效掩膜 & $[96,256]$ & 有效为 1、无效为 0 \\
$X_i$ & 送入模型的物理特征 & $[2,96,256]$ & SSP 通道、地形通道 \\
$V_\ell$ & 第 $\ell$ 个 FNO block 的隐状态 & $[H,N'_z,N'_r]$ & 无量纲特征 \\
$f_i$ & 网络输出的标准化残差 & $[1,96,256]$ & 无量纲 \\
$\hat Y_i$ & 解码后的 TL 预测 & $[96,256]$ & dB \\
\bottomrule
\end{tabularx}
\end{table}

代码采用 PyTorch 的 NCHW 顺序，即一个 batch 张量写成
$[B,C,N_z,N_r]$：$B$ 是 batch size，$C$ 是通道数，最后两维依次是深度和距离。
这一定义贯穿数据预处理、FNO、损失计算和推理，不在中途交换深度与距离轴。
例如 $X_i[1,25,100]$ 表示第 $i$ 个环境中，地形通道在第 25 个接收深度、
第 100 个距离位置上的归一化数值；由于地形只随距离变化，同一距离列的该值相同。
符号 $B$ 只表示一次并行送入网络的样本数，不是海底边界 $b_i(r)$。

\subsection{从一维环境量到二维输入张量}

\subsubsection{SSP 插值、标准化与距离复制}

原始 SSP $c_i(z_j)$ 先用一维线性插值映射到 96 个接收深度 $z_p$：
\begin{equation}
  c^{\mathrm{grid}}_{i,p}=\operatorname{Interp}
  \big(\{z_j,c_i(z_j)\}_{j=1}^{41},z_p\big),
  \qquad c^{\mathrm{grid}}_i\in\mathbb{R}^{96}.
\end{equation}
因为本任务使用距离无关 SSP，同一深度声速沿 256 个距离点复制。声速通道为
\begin{equation}
  C_{i,1,p,q}=\frac{c^{\mathrm{grid}}_{i,p}-1500}{50},\qquad
  C_i\in\mathbb{R}^{1\times96\times256}.
\end{equation}
减去 1500 m/s 并除以 50 m/s 只是稳定数值尺度，不改变 SSP 的垂向形态。

\subsubsection{地形插值、标准化与深度复制}

每条低维 GEBCO 剖面先线性插值到输出的 256 个距离位置。对样本 $i$，地形通道为
\begin{equation}
  B_{i,1,p,q}=\frac{b_i(r_q)}{2000},\qquad
  B_i\in\mathbb{R}^{1\times96\times256}.
\end{equation}
同一距离处的海底深度沿接收深度轴复制，使任意网格点都能访问该距离处的边界位置。
物理输入最终按通道拼接为
\begin{equation}
  X_i=\operatorname{Concat}_{C}(C_i,B_i)
  \in\mathbb{R}^{2\times96\times256}.
\end{equation}
最终模型不使用 Hankel 人工特征；输入信息仅来自 SSP、地形及随后加入的坐标。

\subsubsection{坐标通道与升维层}

FNO 内部生成归一化坐标
\begin{equation}
  Z_{p,q}=\frac{p-1}{N_z-1},\qquad R_{p,q}=\frac{q-1}{N_r-1},
\end{equation}
并把 $[X_i,Z,R]$ 拼成 $4\times96\times256$。逐点 $1\times1$ 卷积把每个
位置的 4 维向量映射到 $H=32$ 维隐藏表示：
\begin{equation}
  V_0=\operatorname{GELU}\big(P[X_i,Z,R]+p_0\big),\qquad
  V_0\in\mathbb{R}^{32\times96\times256},\quad P\in\mathbb{R}^{32\times4}.
\end{equation}
该“升维”不混合相邻网格，只在每个 $(p,q)$ 位置完成通道组合。

\subsubsection{端到端模块数据契约}

表~\ref{tab:tensor-contract} 把公式对应到实际程序接口。方括号均按
\texttt{[batch, channel, depth, range]} 读取；“不改变空间尺寸”表示模块只改变通道内容，
不会删除某个接收深度或距离采样点。

\begin{table}[H]
\centering
\caption{最终 SeaBAR-FNO 的逐模块输入、输出和作用}
\label{tab:tensor-contract}
\begin{tabularx}{\textwidth}{L{2.7cm}L{3.1cm}L{3.1cm}Y}
\toprule
模块 & 输入维度 & 输出维度 & 具体作用 \\
\midrule
物理特征组装 & SSP $[B,41]$；地形 $[B,256]$ & $[B,2,96,256]$ &
插值、标准化，并分别沿距离或深度复制 \\
坐标拼接 & $[B,2,96,256]$ & $[B,4,96,256]$ &
加入深度坐标 $Z$ 和距离坐标 $R$，让算子区分绝对位置 \\
$1\times1$ 升维 & $[B,4,96,256]$ & $[B,32,96,256]$ &
每个网格点独立将 4 个输入量映射为 32 个隐特征 \\
复制填充 & $[B,32,96,256]$ & $[B,32,104,272]$ &
减轻非周期边界直接参与 FFT 所产生的绕回伪影 \\
FNO block $\times4$ & $[B,32,104,272]$ & $[B,32,104,272]$ &
并行计算全局谱路径和逐点局部路径，再归一化、跳连和激活 \\
裁剪 & $[B,32,104,272]$ & $[B,32,96,256]$ &
删除仅服务于计算的 8 行、16 列复制区域 \\
投影头 & $[B,32,96,256]$ & $[B,1,96,256]$ &
逐点执行 $32\rightarrow128\rightarrow1$，输出无量纲残差 $f$ \\
dB 解码 & $f:[B,1,96,256]$ & $\hat Y:[B,96,256]$ &
去掉单通道维，以训练均值场和尺度恢复 TL 单位 \\
\bottomrule
\end{tabularx}
\end{table}

SeaBAR-FNO 可分为五个可独立说明的模块：
\begin{enumerate}
  \item \textbf{双环境上下文编码（Dual-Environment Context Encoder）}：分别把 SSP 和地形
  扩展成二维通道，再与深度、距离坐标融合，使每个预测位置同时知道水体条件、边界条件和位置；
  \item \textbf{各向异性谱传播核（Anisotropic Spectral Propagation Core）}：采用
  $(K_z,K_r)=(16,48)$，在更长的距离轴保留更多 Fourier 模态；
  \item \textbf{全局--局部残差融合（Global--Local Residual Fusion）}：谱路径描述全域耦合，
  $1\times1$ 局部路径保留未截断信息，层间跳连稳定四层堆叠；
  \item \textbf{非周期边界保护（Boundary-Safe Padding）}：复制真实边缘再做 FFT，减轻周期
  绕回。消融显示其收益较小但方向一致，因此视为稳健性增强而非主要精度来源；
  \item \textbf{训练域残差解码（Train-Only Residual TL Decoder）}：网络只拟合训练均值场之上
  的标准化变化，所有统计量均由训练集冻结，避免验证或测试信息进入模型。
\end{enumerate}

\subsection{各向异性谱传播核如何计算}

\subsubsection{复制填充与 Fourier 张量尺寸}

声场在深度和距离边界并不周期。进入第一个谱模块前，在深度末端复制 8 格、距离末端复制
16 格，得到
\begin{equation}
  N'_z=96+8=104,\qquad N'_r=256+16=272,\qquad
  V_0^{\mathrm{pad}}\in\mathbb{R}^{B\times32\times104\times272}.
\end{equation}
对最后两维执行实数二维 FFT 后，因距离轴采用 rFFT 的共轭对称表示，频域张量为
\begin{equation}
  \widetilde V_\ell=\operatorname{rFFT2}(V_\ell)
  \in\mathbb{C}^{B\times32\times104\times137},
\end{equation}
其中 $137=272/2+1$。batch 轴和通道轴都不参与 Fourier 变换。

\subsubsection{各向异性模态截断与复权重混合}

最终模型保留距离方向前 $K_r=48$ 个频率，并在深度方向分别保留顶部和底部各
$K_z=16$ 个频率。对每个保留频率 $(k_z,k_r)$，谱卷积执行从 32 个输入通道到 32 个
输出通道的复数线性变换：
\begin{equation}
  \widetilde S_{\ell,b,o,k_z,k_r}
  =\sum_{h=1}^{32}R^{(\pm)}_{\ell,h,o,k_z,k_r}
  \widetilde V_{\ell,b,h,k_z,k_r},
\end{equation}
其中 $R^{(+)}_\ell,R^{(-)}_\ell\in
\mathbb{C}^{32\times32\times16\times48}$ 分别对应深度频谱的正、负端；未保留模态置零。
逆变换
$S_\ell=\operatorname{irFFT2}(\widetilde S_\ell)$ 恢复到
$[B,32,104,272]$。谱路径的一个输出位置可以利用全域输入，因此适合描述折射和会聚等
长距离耦合；$K_z\ne K_r$ 则显式适配 $96\times256$ 的各向异性网格。

\subsubsection{全局--局部残差融合}

谱路径之外，同一模块还有逐点 $1\times1$ 卷积
$L_\ell=W_\ell V_\ell$，负责不经过频域截断的局部通道混合。完整更新为
\begin{equation}
  V_{\ell+1}=\operatorname{GELU}\!\left(
  \operatorname{GN}\big[S_\ell+L_\ell\big]+V_\ell\right),
  \qquad \ell=0,1,2,3.
\end{equation}
GroupNorm 使用一个 group，在每个样本内部归一化 32 个通道与空间位置；最后的
$+V_\ell$ 是层间残差连接。四个模块均保持 $[B,32,104,272]$ 不变。

\subsection{裁剪、输出头与训练域残差解码}

四个 FNO block 后删除复制区域，恢复为 $[B,32,96,256]$。两个逐点投影层执行
\begin{equation}
  f_i=Q_2\operatorname{GELU}(Q_1V_4+q_1)+q_2,
  \quad Q_1:\,32\rightarrow128,\quad Q_2:\,128\rightarrow1,
\end{equation}
得到 $f_i\in\mathbb{R}^{1\times96\times256}$。$f_i$ 不是 dB 场，而是标准化残差。

仅使用全部训练样本构造一个全局训练均值场
\begin{equation}
  \bar Y(p,q)=\frac{\sum_{i\in\mathcal{D}_{\mathrm{train}}}M_i(p,q)Y_i(p,q)}
  {\sum_{i\in\mathcal{D}_{\mathrm{train}}}M_i(p,q)}.
\end{equation}
将所有训练残差 $Y_i-\bar Y$ 在有效位置拼接，计算单一尺度
\begin{equation}
  s=\max\left\{\operatorname{Std}_{i,p,q:M_i=1}
  [Y_i(p,q)-\bar Y(p,q)],\ 1\ \mathrm{dB}\right\}.
\end{equation}
模型最终输出为
\begin{equation}
  \hat Y_i(p,q)=\bar Y(p,q)+s\,f_i(1,p,q),
  \qquad \hat Y_i\in\mathbb{R}^{96\times256}\ \mathrm{dB}.
\end{equation}
全局均值数组的冻结维度为 $[96,256]$，与网络权重和尺度 $s$ 一同写入 checkpoint；验证
和测试标签从不参与均值、尺度或网络权重估计。均值场本身不区分四条地形，也不根据测试
样本查询某类标签；不同地形造成的差异必须由输入地形通道和 FNO 共同恢复。

\figfull{0.99}{fig04_model_architecture.pdf}{最终 SeaBAR-FNO 的端到端张量维度，以及单个
各向异性 Fourier 模块内部的谱路径、局部路径和 dB 解码。$B$ 表示 batch size。}

\subsection{训练目标、验证选模与单样本推理流程}

训练目标是标准化残差
\begin{equation}
  \widetilde Y_i(p,q)=\frac{Y_i(p,q)-\bar Y(p,q)}{s}.
\end{equation}
只在 Bellhop 有效网格上计算均方损失：
\begin{equation}
  \mathcal{L}(\theta)=
  \frac{\sum_{i\in\mathcal{B}}\sum_{p=1}^{96}\sum_{q=1}^{256}
  M_i(p,q)\left[f_{\theta,i}(p,q)-\widetilde Y_i(p,q)\right]^2}
  {\sum_{i\in\mathcal{B}}\sum_{p=1}^{96}\sum_{q=1}^{256}M_i(p,q)},
\end{equation}
其中 $\mathcal{B}$ 是一个训练 batch。训练使用 AdamW、初始学习率 $10^{-3}$、weight
decay $10^{-5}$、cosine 调度、batch size 16，最多 260 epochs。每个 epoch 后把验证预测
解码回 dB 域计算完整 RMSE；只保存验证 RMSE 最低的 checkpoint，测试集不参与早停或选模。

部署时一个样本依次执行：（1）读取 41 点 SSP 和 256 点地形；（2）形成
$[1,2,96,256]$ 物理输入；（3）网络内部加入坐标并升维；（4）经过 padding 和四个 FNO block；
（5）裁剪并输出 $[1,1,96,256]$ 残差；（6）利用 checkpoint 内的全局均值场和尺度解码为
$[96,256]$ TL。报告的 batch=1 时延从步骤（2）已经形成的 NumPy 特征开始，覆盖输入复制、
步骤（3）--（6）以及输出回 CPU；不包含磁盘读取和原始 SSP 插值。

\section{实验设计与评价方法}

\subsection{对比方法与消融}

\begin{table}[H]
\centering
\caption{对比方法和实验目的}
\begin{tabularx}{\textwidth}{L{3.3cm}L{3.6cm}Y}
\toprule
方法 & 核心设置 & 作用 \\
\midrule
全局训练均值场 & 不使用输入，只输出所有训练场均值 & 主要难度基线；要求测试 RMSE $>4$ dB \\
仅 SSP & 去掉地形通道，其余保持最终结构 & 定量检验显式地形编码是否必要 \\
各向同性谱核 & $(K_z,K_r)=(16,16)$，保留地形 & 检验距离方向增加谱模态的作用 \\
无边界保护 & $(16,48)$，但不执行复制 padding & 检验非周期边界保护的增益 \\
SeaBAR-FNO & 双环境输入、$(16,48)$ 与复制 padding & 最终完整方法 \\
\bottomrule
\end{tabularx}
\end{table}

\subsection{精度指标}

在测试有效网格集合 $\Omega$ 上，误差 $e=\hat y-y$，定义
\begin{align}
  \RMSE &= \sqrt{\frac{1}{|\Omega|}\sum_{(i,z,r)\in\Omega}e_{i,z,r}^2},\\
  \mathrm{MAE} &= \frac{1}{|\Omega|}\sum_{(i,z,r)\in\Omega}|e_{i,z,r}|,\\
  \mathrm{Reduction} &=100\%\times\frac{\RMSE_{\mathrm{mean}}-\RMSE_{\mathrm{model}}}
  {\RMSE_{\mathrm{mean}}}.
\end{align}
除总体指标外，同时报告单样本 RMSE 的中位数、P90 和最差值，以及真实 TL 梯度最高
10\% 网格上的 RMSE，避免总体均值掩盖会聚边缘误差。

\subsection{时延协议与独立复核}

时延测试使用 batch=1，预热 10 次后重复 200 次；GPU 每次显式同步。计时包含 NumPy 到
Torch 输入复制、前向计算、全局均值残差解码和输出回 CPU，不包含磁盘加载。报告中使用 P95
而非单次最小值。训练结束后，独立进程重新计算数据 SHA、重建特征、加载 checkpoint、
只选择 test split 并重新评分和计时。

\section{实验结果}

\subsection{总体结果与相对基线降幅}

正式测试的全局训练均值场 RMSE 为 \GlobalBaseline{} dB，超过 4 dB 难度门槛。最终模型
\BestModel{} 将测试 RMSE 降至 \BestRMSE{} dB，
MAE 为 \BestMAE{} dB，相对全局均值场下降 \RelativeReduction{}\%。

\begin{table}[H]
\centering
\caption{密封测试集总体结果}
\label{tab:results}
\begin{tabular}{lrrrrr}
\toprule
方法 & 参数量/M & RMSE/dB & MAE/dB & P90 样本/dB & A100 P95/ms \\
\midrule
全局训练均值场 & -- & \GlobalBaseline & 3.389 & 6.442 & -- \\
仅 SSP & 6.30 & 4.961 & 3.385 & 6.764 & 4.56 \\
各向同性谱核 & 2.11 & 0.802 & 0.384 & 1.179 & 4.95 \\
无边界保护 & 6.30 & 0.712 & 0.306 & 1.027 & 4.54 \\
SeaBAR-FNO & 6.30 & \BestRMSE & \BestMAE & \BestPninety & \GpuPninetyfive \\
\bottomrule
\end{tabular}
\end{table}

\figfull{0.98}{fig05_main_results.pdf}{全局均值场、仅 SSP 变体和三个含地形算子变体的测试
RMSE，以及相对全局均值场的降幅。}

\subsection{模块消融分析}

仅使用 SSP 时 RMSE 为 4.961 dB，甚至略高于全局均值场；加入地形后，即使采用较紧凑的
各向同性谱核，RMSE 也降至 0.802 dB。这说明在当前多地形任务中，\textbf{显式地形通道是
决定性模块}，而不是可选辅助特征。将距离方向模态从 16 增加至 48 后，无 padding 变体为
0.712 dB，表明各向异性谱设计带来进一步且明确的收益。最后加入复制 padding 后 RMSE 从
0.712 dB 降至 0.702 dB；该增益较小但方向一致，因此报告将其界定为边界稳健性增强。

\subsection{按月份与地形分组结果}

三个月份分组 RMSE 分别为 0.658、0.740 和 0.706 dB；四类地形分组 RMSE 为
0.469--0.864 dB，所有分组均显著低于 2 dB。分组统计使用同一个冻结模型和同一密封测试集，
不为某个月份或地形单独选模。

\figfull{0.93}{fig06_stratified_results.pdf}{SeaBAR-FNO 在 WOA23 3/6/12 月和四类 GEBCO
地形上的密封测试 RMSE；所有子组均通过 2 dB 门槛。}

\subsection{优化过程与典型场}

\figfull{0.78}{fig07_training_curves.pdf}{四个消融变体的验证 RMSE 轨迹。模型选择仅依赖
验证集，红色虚线为 2 dB 门槛；纵轴采用对数尺度。}

\figfull{0.99}{fig08_prediction_examples.pdf}{最终模型在测试集的中位样本和最差样本：
SSP、Bellhop 参考、代理预测、绝对误差以及 1000 m 深度切片。}

最终模型的 P90 单样本 RMSE 为 \BestPninety{} dB，最差单样本 RMSE 为 \BestWorst{} dB。
二维结果表明，主要传播带和地形相关的沿程结构能够被正确恢复；较大误差仍集中在局部
高梯度边缘，而非形成全场系统偏差。

\figfull{0.96}{fig09_error_analysis.pdf}{测试集平均绝对误差、P95 绝对误差空间分布及单样本
RMSE 直方图。}

\subsection{计算时延}

A100 GPU 上最终模型独立复核的完整推理 P95 为 \GpuPninetyfive{} ms；CPU 诊断 P95 为
\CpuPninetyfive{} ms。最终大模型以 A100 为验收平台；若目标机器只使用 CPU，可采用
各向同性谱核变体，其本次 CPU P95 为 66.51 ms、RMSE 为 0.802 dB。25,600 射线 Bellhop
单场中位耗时为 \BellhopMedian{} s。Bellhop 时间只用于说明标签生成成本，不将不同软件
栈的单次运行简单等同为严格加速比。

\figfull{0.78}{fig10_latency.pdf}{三个含地形 FNO 变体的 batch=1 P95 时延与 25,600 射线 Bellhop
单场中位耗时。纵轴采用对数尺度。}

\section{结果分析与适用范围}

\subsection{为什么相对均值场降幅具有说服力}

只报告模型小于 2 dB 并不足以证明代理算子有价值：如果训练均值场本身已经低于 2 dB，
模型可能只是重复一个近似固定场。本项目把全局均值场大于 4 dB 写入自动验收条件，并在
正式标签生成前通过真实地形方位预筛构造有区分度的小任务域。因此，最终结果同时满足
“任务不是平凡的”和“模型能够显著降低误差”两个条件。

\subsection{全局残差统计量的训练边界}

全局均值场和残差尺度只由 288 个训练样本计算；验证和测试标签不参与统计量或网络参数
估计。该均值场对所有地形和 SSP 完全相同，不能按地形类别查询。测试过程只读取输入 SSP、
输入地形与已经冻结的 checkpoint，因此模型必须从地形通道中学习四条剖面的传播差异。

\subsection{结论边界}

本报告可支持的结论限定为：在四条经筛选的巴士海峡 GEBCO 形态、WOA23 3/6/12 月 SSP、
1 kHz、50 m 源深、50 km 和固定底质的仿真域内，模型能够复现当前 Bellhop 标签并满足
精度及时延指标。测试覆盖已知月份家族内的独立扰动，不等同于对未见月份的外推。当前不宣称：
跨海区、未见地形类别、多频率、任意源深、任意底质或海试
条件下仍保持同等精度。若甲方 Bellhop 的射线、束宽、边界实现或输出定义不同，应使用
相同数据契约重新生成少量标签并复核，而不直接假定数值结果完全一致。

\section{结论}

本项目完成了一个从真实公开环境数据、Bellhop 高质量标签到神经算子快速推理的完整闭环：
\begin{enumerate}
  \item 固定 1 kHz、50 km、50 m 源深和二维输出网格，明确了可检验任务；
  \item 基于 GEBCO 2026 与 WOA23 三个代表月份构造了 384 场小而有区分度的受控数据集；
  \item 通过 25,600 射线收敛审计保证当前 reference 的数值稳定性；
  \item 提出 SeaBAR-FNO，将双环境编码、各向异性谱核、全局--局部残差融合、边界保护
  和训练域残差解码组合为一条可复核计算链；
  \item 在全局均值场 \GlobalBaseline{} dB 的非平凡基线上，将测试 RMSE 降至
  \BestRMSE{} dB，并以 \GpuPninetyfive{} ms P95 满足 100 ms 指标；
  \item 提供数据哈希、配置、代码、独立复核与一键复现入口，支持后续在甲方机器上接续开发。
\end{enumerate}

\appendix
\section{关键配置与版本信息}

\begin{longtable}{L{4.4cm}L{9.6cm}}
\toprule
项目 & 冻结值 \\
\midrule
\endfirsthead
\toprule
项目 & 冻结值 \\
\midrule
\endhead
数据配置 & \path{configs/realistic_seasonal_terrain_mvp.yaml} \\
实验配置 & \path{configs/seasonal_campaign.yaml} \\
数据集 ID & \path{bashi_gebco_four_terrain_woa23_m03_m06_m12_v0.7_n384} \\
样本数与划分 & 384；train/validation/test = 288/48/48 \\
输出网格 & 96 depth $\times$ 256 range \\
正式射线数 & 25,600/field \\
数据 SHA-256 & \path{9c90e051851e24082816b435bc529250f64d8a61c8a8f603916071eb0fda0a61} \\
最终模型 & \BestModel \\
代码版本 & Git tag \path{technical-report-v1.5}，具体 SHA 见仓库 release \\
大数据根目录 & \path{/mnt/data/xuangu-fang/ocean-acoustics/projects/ocean-acoustic-surrogate} \\
\bottomrule
\end{longtable}

\section{一键复现与学生接续开发}

\subsection{仓库与目录}

甲方机器建议保持两个兄弟仓库：
\begin{lstlisting}[style=repro]
acoustic-work/
  ocean-acoustic-agent/       # Bellhop backend and task schema
  ocean-acoustic-surrogate/   # dataset, FNO, evaluation and report
\end{lstlisting}

大型标签和 checkpoint 不提交 Git，统一放到：
\begin{lstlisting}[style=repro,language=bash]
export OCEAN_SURROGATE_ROOT=/mnt/data/xuangu-fang/ocean-acoustics/\
projects/ocean-acoustic-surrogate
\end{lstlisting}

\subsection{一键入口}

\begin{lstlisting}[style=repro,language=bash]
cd acoustic-work/ocean-acoustic-surrogate

# 1. Only install/sync and test; Bellhop is not executed.
REPRO_MODE=check bash scripts/reproduce_mvp.sh

# 2. Train SeaBAR-FNO from an existing n384 dataset.
REPRO_MODE=train bash scripts/reproduce_mvp.sh

# 3. Full pipeline: solver check, convergence, labels, training, verification.
REPRO_MODE=full bash scripts/reproduce_mvp.sh

# 4. Verify an existing checkpoint in a fresh process.
export REPRO_RUN_DIR="$OCEAN_SURROGATE_ROOT/runs/<run_id>"
REPRO_MODE=verify bash scripts/reproduce_mvp.sh
\end{lstlisting}

脚本默认使用本报告的 seasonal config、384 样本和最终 SeaBAR-FNO 实验；也可通过
\path{REPRO_CONFIG}、\path{REPRO_CAMPAIGN}、\path{REPRO_EXPERIMENT} 和
\path{REPRO_SAMPLES} 显式覆盖。生成器以每个样本的 \path{sample.npz + metadata.json}
作为断点，重复执行会跳过已完成样本。

\subsection{等价分步命令}

\begin{lstlisting}[style=repro,language=bash]
uv sync --locked
uv run ruff check .
uv run pytest -q

uv run ocean-acoustic-surrogate pilot \
  configs/realistic_seasonal_terrain_mvp.yaml --samples 8
uv run ocean-acoustic-surrogate generate \
  configs/realistic_seasonal_terrain_mvp.yaml --samples 384
uv run ocean-acoustic-surrogate profile \
  configs/realistic_seasonal_terrain_mvp.yaml --samples 384
uv run ocean-acoustic-surrogate campaign \
  configs/realistic_seasonal_terrain_mvp.yaml configs/seasonal_campaign.yaml \
  --samples 384 --only seabar_fno
uv run ocean-acoustic-surrogate verify \
  configs/realistic_seasonal_terrain_mvp.yaml "$REPRO_RUN_DIR" \
  --samples 384 --device auto
\end{lstlisting}

\section{核心代码位置与参考实现}

\begin{table}[H]
\centering
\caption{后续开发的主要代码入口}
\begin{tabularx}{\textwidth}{L{5.2cm}Y}
\toprule
文件 & 职责 \\
\midrule
\path{configs/realistic_seasonal_terrain_mvp.yaml} & 环境、月份、地形、SSP、Bellhop、split 和验收门槛 \\
\path{configs/seasonal_campaign.yaml} & SeaBAR-FNO、训练超参数和模块消融实验 \\
\path{src/ocean_acoustic_surrogate/ssp.py} & Latin hypercube 和平滑 SSP 构造 \\
\path{src/ocean_acoustic_surrogate/dataset.py} & Bellhop task、收敛、断点生成和数据打包 \\
\path{src/ocean_acoustic_surrogate/features.py} & SSP/地形输入和全局残差变换 \\
\path{src/ocean_acoustic_surrogate/models/fno.py} & 谱卷积、padding、残差 block 和逐点输出头 \\
\path{src/ocean_acoustic_surrogate/training.py} & 训练、验证选模、测试和时延 \\
\path{src/ocean_acoustic_surrogate/verification.py} & 新进程 checkpoint 重载和密封测试复核 \\
\path{scripts/reproduce_mvp.sh} & 一键复现入口 \\
\bottomrule
\end{tabularx}
\end{table}

\subsection{生成一个 Bellhop reference}

\begin{lstlisting}[style=repro,language=Python]
from ocean_acoustic_agent import run_simulation
from ocean_acoustic_surrogate.config import MVPConfig
from ocean_acoustic_surrogate.dataset import build_task
from ocean_acoustic_surrogate.ssp import build_ssp_records

cfg = MVPConfig.from_yaml("configs/realistic_seasonal_terrain_mvp.yaml")
record = build_ssp_records(
    cfg.ssp_family, 1, cfg.contract.seed,
    template_cycle_stride=len(cfg.contract.bathymetry.profiles),
)[0]
task = build_task(
    cfg, record, cfg.contract.reference_num_rays,
    task_id="reference_smoke_00000",
)
result = run_simulation(task)
if result.status == "failed" or result.tl_field is None:
    raise RuntimeError(result.error_message)
print(result.tl_field["tl_db"].shape)
\end{lstlisting}

正式生产应调用 \path{generate} 命令，而不是手写循环；正式生成器额外保存 case、配置哈希、
地形类别、split、耗时、有效掩膜和失败记录。

\subsection{调用训练和独立复核}

\begin{lstlisting}[style=repro,language=Python]
from ocean_acoustic_surrogate.config import MVPConfig, load_campaign
from ocean_acoustic_surrogate.training import run_experiment
from ocean_acoustic_surrogate.verification import verify_run

cfg = MVPConfig.from_yaml("configs/realistic_seasonal_terrain_mvp.yaml")
campaign = load_campaign("configs/seasonal_campaign.yaml")
experiment = next(
    item for item in campaign["experiments"]
    if item["id"] == "seabar_fno"
)
dataset = cfg.dataset_root / "n384/dataset.npz"
run_dir = run_experiment(
    cfg, dataset, experiment, campaign["training_defaults"]
)
verification = verify_run(cfg, dataset, run_dir, device_name="auto")
print(verification)
\end{lstlisting}

\subsection{复现通过判据}

一次完整复现至少保留 \path{dataset.npz}、\path{manifest.json}、
\path{convergence_report.json}、\path{model.pt}、\path{metrics.json}、
\path{history.json}、\path{predictions.npz} 和
\path{independent_verification_<device>.json}。验收同时检查
\begin{equation}
  \RMSE_{\mathrm{test}}\le2\ \mathrm{dB},\quad
  \mathrm{MAE}_{\mathrm{test}}\le2\ \mathrm{dB},\quad
  \mathrm{P95}_{\mathrm{latency}}\le100\ \mathrm{ms},\quad
  \RMSE_{\mathrm{global\ mean}}>4\ \mathrm{dB}.
\end{equation}
若复现不一致，依次检查数据 SHA、两个仓库 Git SHA、配置哈希、Bellhop field mode 与射线
数、split、模型 ID 和 CUDA/FFT 后端；不应通过修改测试划分或放宽误差阈值解决。

\section{参考资料}

\begin{enumerate}[label={[\arabic*]}]
  \item Z. Li et al., ``Fourier Neural Operator for Parametric Partial Differential Equations,''
  ICLR, 2021.
  \item M. B. Porter, \textit{The BELLHOP Manual and User's Guide}, 2011.
  \item F. B. Jensen, W. A. Kuperman, M. B. Porter, and H. Schmidt,
  \textit{Computational Ocean Acoustics}, 2nd ed., Springer, 2011.
  \item NOAA NCEI, \textit{World Ocean Atlas 2023},
  \url{https://www.ncei.noaa.gov/access/world-ocean-atlas-2023/}.
  \item GEBCO Compilation Group, \textit{GEBCO Gridded Bathymetry Data},
  \url{https://www.gebco.net/data-products/gridded-bathymetry-data}.
  \item IOC, SCOR and IAPSO, \textit{The International Thermodynamic Equation of Seawater
  -- 2010: Calculation and Use of Thermodynamic Properties}, 2010.
\end{enumerate}

\vfill
\begin{center}
\color{gray}\small —— 报告结束 ——
\end{center}

\end{document}
